\section{No-Other-Solutions Standard}
\label{sec:pr35-no-other-solutions-standard}

For each sector $s$, let \(\mathfrak A_s\) be the complete admissible
candidate set, let \(\sim_s\) identify gauge, basis, phase, coordinate, and
unit conventions, and let \(F_s\) collect all preregistered mathematical
constraints.  The uniqueness claim requires
\begin{equation}
\boxed{
\{x\in\mathfrak A_s:F_s(x)=0\}/\!\sim_s
=
\{[x_{\rm PR}]\}.}
\label{eq:pr35-sector-uniqueness-standard}
\end{equation}
Verifying only \(F_s(x_{\rm PR})=0\) establishes existence, not uniqueness. The proof must also show local isolation after unphysical null directions are removed and exhaust every disconnected admissible branch. Finite discrete candidates are treated by exhaustive enumeration. Polynomial and rational constraints are solved by exact elimination or complete analytic parameterization. Continuous numerical roots require interval isolation and exclusion of the complement. Operator claims require self-adjoint domains and certified spectral enclosures. Radiative corrections require a complete invariant-operator basis and an explicit power-counting rule.

Post-generation diagnostic agreement is not an admissibility constraint. If \(\mathcal E_s\) denotes external data and \(\mathcal C_s\) the construction map, the firewall condition is 
\begin{equation}
\mathcal E_s\cap\operatorname{Dom}(\mathcal C_s)=\varnothing .
\end{equation}
No residual or standard-deviation score may select a branch, coefficient, truncation order, numerical resolution, or acceptance tolerance. The final global theorem is permitted only when every sector contains one surviving equivalence class or is explicitly excluded from the theorem scope.
